\newpage
\section{Methodology\label{sec:methods}}

An analysis of the growth of streetcar systems in the United States requires data on the length of systems in each state as well as their historical populations. The former was obtained from the McGraw-Hill Directories between 1894 and 1932. The analysis presented herein runs til 1926, as the later years (1927-1932) presented yet another data formatting challenge that could not be overcome within the time allotted. The population data was obtained from the National Historical Geographic Information System, hereafter NHGIS, records of the 1910 and 1920 United States national census \citep{nhgis}. Once the data was extracted, a three-parameter logistic function along with regression was applied to describe an S-curve modelling the growth of streetcar systems and evaluate its goodness of fit. In addition, a comparison was run between human-language and machine-reading processing to compare the advantages and disadvantages of sorting data via code.

This section is split into four subsections. The first deals with extraction of the data from the McGraw Directories, using flowcharts to explain the logic behind Python coding presented in abbreviated, annotated form in Appendix D and in full in Appendix E. The second explains the process and equations used to generate the S-curve, of which coding is provided for in Appendix E. The third is a brief section introducing the statistical package used to evaluate goodness of fit and other relevant statistical measures, and the fourth is a section comparing human-language and machine-reading processing. 

\subsection{Extracting the data}
Although the McGraw-Hill Directories offer an in-depth look at streetcar systems in the United States, their formatting was inconsistent between the years. Hence, this section is further subdivided according to the distinct formatting periods in the Directories. 

The Directories were rendered into text files using version 12.1.14 of the software FineReader OCR Pro \citep{ocr}. Texts were preprocessed by adding 'START OF RECORDS' when records began and 'END OF RECORDS' when records concluded, skipping over sections displaying advertising and records of business owners. Where relevant, an 'INDEX TO STREET RAILWAY COMPANIES' and 'END OF COMPANIES' were added to the start and end of the indexes.

Note that the term 'miles of track' refers to single track, including main track, second track and sidings owned by the company. The code attempts to avoid miles of leased, operated, and under-construction track, although this was not always successful, and results had to be manually checked for double-counted tracks.

\subsubsection{Directories from 1894 to 1902}

Between 1894 and 1902, the McGraw Directories were sorted alphabetically by city, with same-name cities organised by state. Using the index of companies (or the 1898 index, in the case of 1894 which provided no index), matching was done on a city-company-state basis; that is, if the recorded city and company matched the index's city and company, the state was recorded. 

In extracting the index of companies, the read\_file function \citep{readfile} and the module pandas \citep{mckinney2010data} were used for interpreting the text file. The flowchart is presented in Figure 3.1, and relies on the presence of 'Co.', 'Ry.' and 'R. R', the latter two short for railway and railroad respectively, in company names. The check for capitalisation exists to avoid strings such as 'DIGITIZED BY GOOGLE' and 'AMERICAN STREET RAILWAY INVESTMENTS' After strings were cleaned and the csv file exported, a brief manual edit was performed to ensure all entries in the second column, 'Place', were formatted as 'CITY, STATE'. 

\begin{figure}[htp]
    \centering
    \includegraphics[width=\textwidth]{figures/indexex-1.png}
    \caption{Flowchart for extraction of company index}
\end{figure}

In finding the cities, two separate functions were designed: one to extract it from the text file, presented in Figure 3.2, and one to match it to the cleaned value in the index, presented in Figure 3.3. The modules used were pandas, numpy \citep{2020NumPy-Array} and fuzzywuzzy \citep{fuzzywuzzy}. 

\begin{figure}[htp]
    \centering
    \includegraphics[width=\textwidth]{figures/get_cities_to14-1.png}
    \caption{Flowchart for function used to identify potential cities} 
\end{figure}

\begin{figure}[htp]
    \centering
    \includegraphics[width=\textwidth]{figures/clean_cities_to14-1.png}
    \caption{Flowchart for function used to confirm cities.}
\end{figure}

To extract the miles from the data, a function was created to search for the phrase 'Plant and Equipment' and extract the first number after the phrase. It relied on a consistent formatting for the length of a streetcar system being presented as 'Plant and Equipment--Miles of track (mode), (number)', with additional instructions to avoid leased track but pick up secondary track. However, some miles were represented in a table, which the OCR was unable to interpret, and these had to be manually checked by hand. The flowchart for this function is presented in Figure 3.1.4, and the modules used were pandas, numpy, fuzzywuzzy and re \citep{van1995python}. 

\begin{figure}[htp]
    \centering
    \includegraphics[width=\textwidth]{figures/get_milesto14-1.png}
    \caption{Flowchart for function used to extract miles.}
\end{figure}

For companies, checks were first run on the dataframe for specific phrases, in order to distinguish between company listings and subcompany listings. Subcompany listings were more frequent in larger cities such as Boston and New York, where the Directories listed a number of smaller companies under larger organisations such as the Massachusetts Electric Companies. Although these smaller companies did not have their track listed separately, they were still included in the index, making it unwise to apply the index directly. Companies and subcompanies were distinguished by the presence of 'Capital Stock', 'Funded Debt', 'Equipment' or 'Population' listed near the company. After this PhraseCheck was performed, the function reasoning in Figure 3.1.5 could proceed. The function used the pandas, re and fuzzywuzzy modules.

While the minimum similarity limit given in Figure 3.1.5 is 90\%, the company extraction function was run several times with decreasing limits, up to an 80\% similarity. This was both to capture companies who had not been accurately captured by OCR and to avoid cases like 'Second Avenue Railway' matching with 'Ninth Avenue Railway'.

\begin{figure}[htp]
    \centering
    \includegraphics[width=\textwidth]{figures/get_companies_to14-1.png}
    \caption{Flowchart for function used to extract companies.}
\end{figure}

Finally, once the city and company were acquired, the state was found by matching the two in the index. The dataframe was sliced to only the lines containing the miles of track, and the data was extracted as a csv file. 

\subsection{Directories from 1903 to 1914}
Between 1903 and 1914, the McGraw Directories were reorganised to alphabetically by state. This did not affect the logic of the overall code, which still processed the dataframe through index extraction, city identification, company identification and state identification, though a line of code used to split city from state in the code represented by the flowchart in Figure 3.1.2 was removed.

Of note during this period was an increase in miles of track entries where companies and their track was leased and counted under a larger company, but also listed in its own entry. It was difficult to separate these cases, as the Directories appeared to distinguish between 'all capital stock owned by another company' and 'leased to another company' without clear reason. While the code caught the ones listed with 'leased' or 'included' in the Plant and Equipment line, other entries had to later be manually removed.

Another change in this period was from 1911 onward, where 'Plant and Equipment' changed to 'Track and Equipment'. Other than replacing the phrase 'Plant and Equipment' with 'Track and Equipment' in the function to extract miles, no change was made. 

\subsection{Directories from 1917 to 1926}

From 1917 onwards, the McGraw-Hill Directories underwent another significant change. In some years, two copies of the directories were produced; it was decided to use only one of the copies. The August edition was used for 1917, 1919, 1920, 1921 and 1922, and the February edition for 1918 and 1923, as in some cases one copy was not available. 

In these editions, many references to operations, balance sheets, passenger revenue and other numbers of interest to potential stock investors disappeared, possibly being published in the accompanying "McGraw Electrical Trade Directory", and abbreviations began to be used more frequently. This made the fuzzywuzzy module unable to reliably match companies using the index without dropping to what were deemed as unacceptably low minimums, and so new functions had to be devised.

Owing to the new and more compressed format, the OCR began returning inconvenient line breaks, and so before running the code the data underwent preprocessing to remove most line breaks in Notepad++. Line breaks before and after majority-capitalised strings with numbers were kept, as these strings roughly corresponded to the cities, while other line breaks were removed to enable better searching for companies. 

Due to being unable to use the index, the module 'us' \citep{usmodule}, which provided a list of American states and territories, was used to find states instead. It provided a list of states against which fuzzywuzzy could be used to match string and state together.

Most cities in the 1917-1926 Directories were formatted as 'CITY, POPULATION', which the function to record cities made use of. By splitting the string, it was possible to acquire the city on its own, although punctuation was sometimes present and was removed after. A notable filter necessary in the function was for parts of the string 'POWER STA. EQUIP', which often came with numbers and threw the generated database into disarray. The flowchart in Figure 3.3.1 presents the logic behind the new function, with the modules pandas, fuzzywuzzy, numpy and re being used. 

\begin{figure}[htp]
    \centering
    \includegraphics[width=\textwidth]{figures/get_cities_from17-1.png}
    \caption{Flowchart for function used to extract cities post-1917.}
\end{figure}

Companies within the 1917-1926 McGraw Directories tended to come in the form '(Number) — (Company) — remaining string', such as "596—North Jersey Rapid Transit Co.— Gen. Office, Hobokus. [sic] (Connects Paterson. Fairlawn Glen..." The code for companies made use of the numbers and emdash to write directly from the McGraw Directories, with string cleaning performed post-extraction. The modules used were pandas, fuzzywuzzy, numpy and re.

\begin{figure}[htp]
    \centering
    \includegraphics[width=\textwidth]{figures/get_companies_from17-1.png}
    \caption{Flowchart for function used to extract companies post-1917.}
\end{figure}

Although not mentioned in fuzzywuzzy's documentation, the partial matching algorithm loses reliability once strings pass 200 characters. This increased the difficulty of obtaining the length of track, as the sentence containing this information was often located near the end of long strings. As a workaround, four phrases to potentially split the string were identified, allowing fuzzywuzzy to work on the shortened strings. The other modules used were re and pandas. 

\begin{figure}[htp]
    \centering
    \includegraphics[width=\textwidth]{figures/get_companies_from17-1.png}
    \caption{Flowchart for function used to extract miles post-1917.}
\end{figure}

\subsection{Generating the S-curve}
To model the data, \autoref{eq:scurve}, a three-parameter logistic equation, was used:
\begin{equation}
\label{eq:scurve}
    S(t) = \frac{K}{1+e^{-b(t-t_i)}}
\end{equation}

\noindent
where $S(t)$ = miles of track,\\
$K$ = variable estimated from regression, corresponding to the asymptotic maximum of the system\\
$t$ = time (years)\\
$t_i$ = year in which 0.5 $S_{max}$ is achieved, and\\
$b$ = an estimated coefficient describing the rate of growth\\

The coefficients were estimated by regression, in the form of \autoref{eq:regression}:

\begin{equation}
    Y = bX + c
\label{eq:regression}
\end{equation}

\noindent where $Y = ln(\frac{Track}{K - Track})$, \\
$X=$ year\\

For systems that had entered a decline by 1926, the K-value was found by a short iteration of values between the combined length of track at the height of the system and the number 10 above that, and choosing the K-value with the highest correlation coefficient, $R^2$. This was necessary as the K-value represents an asymptote, not merely the extent of the system, and +10 provided a reasonable asymptote close to the maximum while still allowing for the correlation coefficient to adjust the asymptotic maximum. Once the K-value was identified and a list of Y-values produced from equation 3.2, the curve fitting function np.polyfit() was used to identify $b$ and $c$, through which $t_i$, the inflexion year could be obtained.

Where systems had not entered a notable decline by 1926, the iteration for K was increased to between the highest combined length of track and 1.5x the value, again choosing the K-value with the highest correlation coefficient to find $b$ and $c$. As historically streetcars had entered decline by the 1920s and would continue declining into the 1930s, 1.5x was considered a suitable multiplier while limiting the equation's potential for extreme and historically accurate growth. Again, the function np.polyfit was used to identify the rate of growth and the inflexion year. The curves and the calculated S-curve was then plotted using the module matplotlib \citep{hunter2007matplotlib}. 


\subsection{The Maturity and Decline Phase}
While generating the S-curve, it was discovered that for areas which had entered a period of decline it was necessary to split the data into two; one to cover the period of growth, and one to cover the period of decline. Without separating the periods, the S-curve became extremely inaccurate and the $R^2$ value drops accordingly. Figure 3.5.1 and Table 3.5.1 represent the S-curve for the state of Massachusetts, a state with a noticeable growth and decline period, without breaking the data into growth/decline periods, while Figure 3.5.2 and Table 3.5.2 display the much closer-fit S-curves once the dataset is split.

\begin{figure}
    \centering
    \begin{minipage}{0.45\textwidth}
        \centering
        \includegraphics[width=0.9\textwidth]{figures/MassachusettsNoBreak.png} % first figure itself
        \caption{Massachusetts S-curve without separation}
    \end{minipage}\hfill
    \begin{minipage}{0.45\textwidth}
        \centering
        \includegraphics[width=0.9\textwidth]{figures/Massachusetts.png} % second figure itself
        \caption{Massachusetts S-curve with separation}
    \end{minipage}
\end{figure}

\begin{table}[h!]
\centering
\begin{tabular}{lrrrrr}
\hline
State & Sample Size & K-value & RSQ-adjusted & Standard error & P-value \\
\hline
Massachusetts & 26 & 6095 & -0.00162 & 0.554  & 0.337  
\end{tabular}
\caption{Massachusetts statistics without separation}
\end{table}

\begin{table}[h!]
\centering
\begin{tabular}{lrrrrr}
\hline
State & Sample Size & K-value & RSQ-adjusted & Standard error & P-value \\
\hline
Massachusetts & 15 & 4073 & 0.833 & 0.722 & 1.27E-06 \\
\hline 
Massachusetts & 14 & 4073 & 0.693 & 0.911 & 4.78E-04
\end{tabular}
\caption{Massachusetts statistics with separation, regressing to zero}
\end{table}

Separation of years into growth and decline periods were done for states that had reached a maximum before 1923. For these states, data was split between the years up to and including the highest track value, and the years including and after the highest track value, counting the year with the highest track value twice in order for the selected K-value to apply to both periods as an appropriate asymptote. 

Three separate versions of the decline S-curve were calculated. The first version was done with only the values in the McGraw-Hill directories, letting the lower limit be zero in accordance with historical records. The second and third versions both featured an extra entry to test for a closer S-curve model to the values up to 1926. In the second model, an additional entry of [1940, one third of maximum track] was added, and in the third model an additional entry of [1940, 90\% of lowest record during decline phase] was tested. No changes were made to the equation or data.

\subsection{Statistical analysis}
Statistical analysis was performed using the Python module statsmodel \citep{seabold2010statsmodels} for Ordinary Least Squares. The following parameters were extracted from the model for each area of interest: intercept, slope, adjusted $R^2$ value, p-value, and three measures of standard error for the regression, intercept and slope respectively. 

All results are recorded in Appendixes A and B.

\subsection{Data Validation}
Data validation of the code-extracted results was performed via comparison to hand-extracted values done as unpublished classwork by the \citeauthor{civlclass} at the University of Sydney. Comparison graphs were made for each state done by the classes, as well as an aggregated all-states comparison. The data validation was performed for the Directories up to 1920, as the 1921-1926 Directories were not located until later. 

